\chapter{Conclusão}
Este trabalho teve como objetivo desenvolver uma abordagem automatizada para a resolução de jogos de adivinhação de palavras em português, utilizando técnicas de processamento de linguagem natural. A partir do desenvolvimento, implementação e experimentação da solução proposta, é possível afirmar que todos os objetivos específicos estabelecidos foram plenamente alcançados.

% A revisão dos trabalhos relacionados evidenciou que há múltiplas abordagens propícias para este tipo de problema, enfatizando o potencial dos modelos de aprendizado de máquina na captura de relações semânticas entre termos. Com base nesse embasamento teórico, foram selecionados três modelos amplamente utilizados — Word2Vec, GloVe e FastText —, todos treinados em corpus da língua portuguesa, a fim de avaliar o desempenho de cada um no desafio proposto.

% Os resultados obtidos demonstraram que o sistema desenvolvido, implementado em Python e integrado ao jogo por meio do Playwright, foi capaz de encontrar a palavra secreta em uma proporção significativa das execuções realizadas, comprovando a eficácia da abordagem. Entre os modelos avaliados, o GloVe apresentou o melhor desempenho, destacando-se por alcançar a maior taxa de acertos.

% Portanto, constata-se que a utilização de embeddings semânticos para automatizar a resolução de jogos linguísticos não é apenas viável, mas também eficaz, reforçando o potencial dessas técnicas em aplicações voltadas ao entretenimento, à educação e ao aprimoramento de modelos de linguagem em português.

Foi desenvolvido um sistema completo para a resolução automática de partidas do jogo Contexto, incluindo a estruturação da classe ContextoSolver, a integração com o Playwright e a criação de módulos auxiliares. Em seguida, uma heurística de seleção de palavras foi implementada, combinando similaridade semântica, exploração inicial e um modo de otimização com Regressão Ridge, abordagem responsável por guiar o processo de busca e aprimorar dinamicamente as escolhas do algoritmo durante a partida.

Quanto à avaliação da eficiência dessa heurística, foram conduzidos experimentos em mil execuções, registrando taxa de acertos e quantidade de tentativas. Os resultados demonstraram que a heurística é funcional, consistente e capaz de resolver uma parcela significativa das partidas.

Além disso, foi realizada uma investigação estruturada de diferentes modelos de word embeddings incluindo o estudo de suas características e sua aplicação prática na resolução automática. Tal exploração permitiu compreender suas diferenças, capacidades e limitações no contexto do projeto. O sistema desenvolvido foi capaz de comparar o desempenho desses modelos sob a mesma heurística, garantindo uma análise justa e padronizada.

Com isso, o projeto cumpriu seus objetivos específicos, tais como a criação de um sistema funcional de automação, o desenvolvimento de uma heurística eficaz de seleção de palavras, a avaliação de sua eficiência por meio de experimentos sistemáticos e a comparação estruturada entre diferentes modelos de word embeddings, destacando o GloVe como o mais adequado à tarefa.

\section{Trabalhos Futuros}

Como desdobramento deste estudo, é possível diversas oportunidades de aprimoramento e expansão da solução proposta. Uma primeira vertente consiste na otimização da estratégia de seleção de palavras através de algoritmos de Aprendizado por Reforço, como o Q-Learning. Conforme explorado na revisão bibliográfica, essa abordagem permitiria que o agente aprendesse uma política de escolha ótima baseada nas recompensas (ranking) recebidas a cada tentativa, visando reduzir o número total de interações necessárias para encontrar a palavra secreta.

Outra possibilidade relevante é a implementação de uma estratégia de heurística híbrida ou dinâmica. Um trabalho futuro poderia investigar o uso de um modelo com maior capacidade de generalização para a etapa de exploração inicial, alternando automaticamente para um modelo com maior precisão local ou ajustado para sinônimos na etapa de refinamento, quando o ranking se aproxima do alvo. Adicionalmente, técnicas de Ensemble, que combinam a votação de múltiplos modelos para decidir a próxima jogada, poderiam mitigar as fraquezas individuais de cada arquitetura.

Recomenda-se também a investigação de abordagens baseadas em \textit{Graph Embeddings} e grafos de conhecimento. Diferente dos modelos vetoriais tradicionais, que dependem puramente de distâncias cosseno em um espaço denso, modelar o vocabulário como um grafo, onde os vértices são palavras e as arestas representam o peso da proximidade semântica, permitiria explorar algoritmos de caminho mínimo ou detecção de comunidades. Isso poderia solucionar casos onde a similaridade vetorial falha em capturar relações diretas, como hiperônimos ou termos técnicos específicos.

Por fim, propõe-se a evolução da ferramenta para suportar um formato de torneio de agentes autônomos, inspirado na dinâmica de competições de robótica. Nessa estrutura, o sistema desenvolvido atuaria como uma "arena" padronizada, aproveitando sua arquitetura modular. Pesquisadores e estudantes seriam desafiados a submeter seus próprios "robôs de software", compostos por modelos treinados em corpora personalizados e heurísticas ajustadas, para competir na resolução simultânea de desafios. Essa abordagem fomentaria a busca prática por eficiência, premiando soluções que demonstrem maior criatividade na engenharia de dados e na otimização algorítmica.